% =====================================================================
%  1bat-ud2-16.tex  —  HORA 16: Prova escrita de la Unitat 2
%  (sense preàmbul: s'inclou des de main.tex)
% =====================================================================
\horatitol{Sessió 16 — Prova escrita de la Unitat 2}%
{Prova individual: quatre problemes, un per cada bloc treballat (C1–C4).}

\subsection{Instruccions}

\begin{itemize}[leftmargin=1.4em, itemsep=2pt]
  \item \textbf{Durada:} 50 minuts.
  \item \textbf{Material permès:} calculadora científica. No es permet
        l'ús del mòbil ni de cap altre dispositiu connectat.
  \item La prova consta de \textbf{quatre problemes}, un per a cada bloc
        de la unitat (C1–C4), amb el \textbf{mateix pes} cadascun.
  \item Es valora tant el \textbf{resultat} com el \textbf{procediment}:
        mostra sempre els passos, controla el signe i el grau, i
        interpreta el resultat en el context del problema (vegeu la
        rúbrica).
  \item Escriu la resposta a l'espai indicat sota de cada problema.
\end{itemize}

\subsection{Rúbrica}

Cada problema es corregeix amb els mateixos quatre nivells treballats al
taller de la sessió 15:

\begin{center}
\renewcommand{\arraystretch}{1.3}
\begin{tabular}{|c|p{9.7cm}|}
\hline
\textbf{Nivell} & \textbf{Descripció} \\
\hline
\textbf{NA} — No assolit &
  No es resol el problema, o es planteja amb un error greu que en
  desvirtua el sentit. \\
\hline
\textbf{AS} — Assoliment satisfactori &
  El resultat final és correcte, però sense procediment visible ni
  justificació dels passos. \\
\hline
\textbf{AN} — Assoliment notable &
  El procediment és correcte i justificat, amb algun error menor de
  càlcul, de signe o de grau que no l'invalida. \\
\hline
\textbf{AE} — Assoliment excel·lent &
  El procediment és complet i correcte, amb el signe i el grau ben
  controlats, i el resultat s'interpreta correctament en el context
  del problema. \\
\hline
\end{tabular}
\end{center}

\subsection{Exercicis per fer}

\begin{exercicis}[Prova — 50 minuts]
\begin{exlist}

  \item \textbf{(C1 — Llenguatge algebraic)} Un servei de suport
        familiar aplica un criteri de copagament: una quota fixa de
        $20\eur$ més $2\eur$ per cada hora de suport rebuda. Tradueix
        aquest criteri a una expressió algebraica que doni el
        copagament total, $P(x)$, per a $x$ hores de suport.
  \vspace{3cm}

  \item \textbf{(C2 — Valor numèric)} El copagament d'un servei de
        teleassistència es calcula amb $T(x)=\num{0,05}x^2-9x+450$, on
        $x$ és la renda mensual de la persona usuària, en euros.
        Calcula el copagament que correspon a un expedient amb una
        renda mensual de $60\eur$, és a dir, $T(60)$.
  \vspace{3cm}

  \item \textbf{(C3 — Operacions amb polinomis)}
    \begin{enumerate}[label=\alph*), leftmargin=1.6em, topsep=2pt]
      \item Dos serveis comunitaris comparteixen un projecte. El
            pressupost d'un és $A(x)=4x^2+2x-6$ i el de l'altre
            $B(x)=x^2-5x+9$, on $x$ és el nombre de famílies ateses.
            Calcula el pressupost conjunt $A(x)+B(x)$.
      \item Calcula i simplifica: $3\per(2x^2-x+4)-(x^2+3x-7)$.
    \end{enumerate}
  \vspace{3cm}

  \item \textbf{(C4 — Arrels)}
    \begin{enumerate}[label=\alph*), leftmargin=1.6em, topsep=2pt]
      \item Una prestació es modelitza amb $B(x)=-12x+540$, on $x$ són
            els ingressos mensuals, en desenes d'euros. Troba el
            llindar d'exclusió de la prestació i interpreta el
            resultat.
      \item Resol l'equació $x^2-8x+12=0$ aplicant la fórmula general.
    \end{enumerate}
  \vspace{3cm}

\end{exlist}
\end{exercicis}

\bigskip
\noindent\textit{Graella de correcció — ús del professorat.}
\begin{center}
\renewcommand{\arraystretch}{1.4}
\begin{tabular}{|p{5.6cm}|c|c|c|c|}
\hline
\textbf{Criteri} & \textbf{NA} & \textbf{AS} & \textbf{AN} & \textbf{AE} \\
\hline
C1 — Llenguatge algebraic & $\square$ & $\square$ & $\square$ & $\square$ \\
\hline
C2 — Valor numèric & $\square$ & $\square$ & $\square$ & $\square$ \\
\hline
C3 — Operacions amb polinomis & $\square$ & $\square$ & $\square$ & $\square$ \\
\hline
C4 — Arrels & $\square$ & $\square$ & $\square$ & $\square$ \\
\hline
\end{tabular}
\end{center}

% ---------------------------------------------------------------------
\fullsolucions{Sessió 16}
\begin{solucions}
\begin{sollist}
  \item $P(x)=2x+20$ (en euros, per a $x$ hores de suport).
  \item $T(60)=\num{0,05}\per3600-9\per60+450=180-540+450=90$: el
        copagament és $90\eur$.
  \item (a) $A(x)+B(x)=(4+1)x^2+(2-5)x+(-6+9)=5x^2-3x+3$. \quad
        (b) $3\per(2x^2-x+4)-(x^2+3x-7)=6x^2-3x+12-x^2-3x+7
        =5x^2-6x+19$.
  \item (a) $-12x+540=0 \Rightarrow x=45$ (desenes d'euros, és a dir
        $450\eur$ d'ingressos): a partir d'aquest nivell d'ingressos ja
        no hi ha dret a la prestació. \quad
        (b) $\Delta=(-8)^2-4\per1\per12=64-48=16$, $\sqrt{\Delta}=4$;
        $x=\dfrac{8\pm4}{2} \Rightarrow x_1=6$, $x_2=2$.
\end{sollist}
\end{solucions}
